\documentclass[11pt]{article}
\usepackage{fontspec}
\setmainfont{Times New Roman}
\setsansfont{Arial}
\setmonofont{Consolas}
\usepackage{amsmath,amssymb,mathtools}
\usepackage{geometry}
\usepackage{microtype}
\newcommand{\mo}{\mathfrak{o}}
\newcommand{\mO}{\mathfrak{O}}
\newcommand{\mT}{\mathfrak{T}}
\newcommand{\mR}{\mathfrak{R}}
\newcommand{\mq}{\mathfrak{q}}
\newcommand{\ma}{\mathfrak{a}}
\newcommand{\mb}{\mathfrak{b}}
\newcommand{\mc}{\mathfrak{c}}
\newcommand{\mpideal}{\mathfrak{p}}
\geometry{a4paper,margin=2.35cm}
\setlength{\parindent}{1.25em}
\setlength{\parskip}{0pt}
\makeatletter
\renewcommand\footnotesize{\@setfontsize\footnotesize{7.7}{8.7}\selectfont}
\makeatother
\emergencystretch=100em
\allowdisplaybreaks
\sloppy
\hbadness=10000
\vbadness=10000
\hfuzz=2pt
\vfuzz=10pt
\raggedbottom

\begin{document}

% Унит: Папер31 Сецтион08 Ентры06 мултиплицатион-ринг qуотиент лоцализатион в001. Дирецт Интерславиц Латин транслатион фром цлеан РА34 лине 439 / сегмент P31-С0099, вит сцан-висибле финал еqуалиты/прооф сентенце ресторед фром принтед паге 103 / ПДФ паге 22.
\setcounter{footnote}{37}

\noindent Те саме размысльеньа важе такоже за мултипликацијско колцо $\mo$. Особливо, ако $\mpideal$ јест просты идеал, тогда $\mo_\mpideal$ имаје само једин просты идеал, различны од нулового и јединичного идеала; заради изоморфизма колц класов остатков колцо $\mo_\mpideal$, како и $\mo$, јест мултипликацијско колцо, и по 1 тогда $\mo_\mpideal$ јест колцо главных идеалов. Елемент, кторы пораждаје просты идеал изведены из $\mpideal$, става простым елементом $p$ из $\mo_\mpideal$. Разширјены идеал либовольного идеала $\mc$ из $\mo$ јест породжены примарноју компонентоју $\mc$, принадлежечоју к $\mpideal$; зато он јест равны $(p)^\alpha$ або јединичному идеалу, согласно тому, входи ли $\mpideal$ в $\mc$ -- и то в $\alpha$-тој степени -- або не входи в $\mc$. Из того још следује: ако разширјеньа двух идеалов $\mb$ и $\mc$ из $\mo$ совпадајут в всих колцах квоциентов $\mo_\mpideal$, тогда $\mb$ и $\mc$ сут идентичне, бо они сут делимы тыми самыми простыми идеалами $\mpideal$ и в тых самых степеньах.

\end{document}

